% MT007 / explanation — canonical editable content.
把“阳性”定义为需要进一步确诊检查。漏掉患者叫假阴性（false negative，FN），把健康者送去进一步检查叫假阳性（false positive，FP）。本题允许一些FP来减少更昂贵的FN，因此不能只追求两种错误加起来最少。


\textbf{方法一在训练时改代价。} 提高阳性样本的权重，或专门提高假阴性的损失。这样漏掉一位患者对目标函数的影响更大，模型会更重视把患者放到阳性一侧。它会重新改变学出的权重和边界。


\textbf{方法二在预测时改阈值。} 若\(p=P(\text{患病}\mid x)\)，原来\(p>0.5\)才判阳性，现在可降低门槛。补充例子：某人的模型概率为0.3，门槛0.5时判阴性，门槛0.2时就进入进一步检查。更多人被判阳性通常会减少漏诊，也会增加健康者的检查数量；0.2只是说明方向，不是本题给定或推荐的真实医疗阈值。


两种方法分别改变模型训练和最终决策，属于两个不同答案。用验证数据选择权重、阈值，并观察召回率\(TP/(TP+FN)\)及误报成本；其中\(TP\)是正确识别的患者数。原题目标是“不漏掉病例”，但有限验证数据上的表现不能保证未知患者中绝无漏诊。


选读：若概率已校准，漏诊代价为\(C_{FN}\)、误报代价为\(C_{FP}\)，判阳性的预期代价是\(C_{FP}(1-p)\)，判阴性是\(C_{FN}p\)。前者更小时选阳性，得到\(p>C_{FP}/(C_{FP}+C_{FN})\)。漏诊越贵，这个阈值越低。


\textbf{出处：} 2020B期中Quiz Q7，原卷第3页（10分）；2023B期中 Q13，原卷第10页（10分）


\textbf{答案依据：} 2020B期中 Q7 随卷／配套参考答案；条件与解释经复核。
